\section{The probability of passing a coverage screen}
\label{v13:la-screen}

Evaluating an agent that recommends inferential procedures requires a clear
distinction between nominal coverage, an argument for coverage and an empirical
qualification rule. Suppose a fixed procedure is tested in $n$ independently
generated replications of a declared mechanism. Let $K$ count requests whose
reported sets cover the truth; unavailable or failed requests remain in the
declared denominator. If the per-request coverage is $p$, then under these
simulation assumptions $K$ has a binomial law.

For a simultaneous family of $m$ cells and total error budget $\alpha$, consider
the one-sided exact lower bound
\[
 L(K,n)=\operatorname{BetaQuantile}
       (\alpha/m;K,n-K+1),\qquad K>0,
\]
with $L(0,n)=0$. A rule requiring $L\geq0.95$ certifies a minimum in this
specific finite experiment. At true $p=0.95$, its probability of passing one
cell is at most $\alpha/m$. This is how a lower-bound certification test behaves
at its boundary, not evidence that the procedure's true coverage is zero.

Before qualification, a designer can compute the smallest integer $k_*$ that
passes and the operating characteristic
\[
 \Pr_p(\hbox{pass})=\Pr_p(K\geq k_*).
\]
The project calculation evaluates this exactly for several replication counts,
family sizes and values of $p$, and independently checks the equivalent
binomial-tail threshold. It draws no random data and does not test a causal
estimator. Its purpose is to make the qualification design interpretable.
\cite{v13_project_exact}

Increasing replications can distinguish a genuine coverage margin above 95\%
more sharply. It cannot make a level-controlled test have high power at the
exact boundary. This does not justify lowering the required minimum after
seeing results. The analyst must state what guarantee is theoretically justified,
what empirical screen is being applied and what its power is. A method may
have a plausible asymptotic argument while failing this finite screen; that
should be reported as a qualification shortfall rather than silently recoded
as success.

An agent could otherwise exploit this distinction in two bad ways. It might
declare a method invalid merely because a low-powered boundary screen failed,
or widen intervals until they pass while ignoring decision usefulness. The
evaluation should penalize both. Coverage, finite availability and useful width
are separate dimensions. A whole-line interval covers every finite truth but
may support no action. Bank width and loss thresholds require actual decision
owners; synthetic thresholds should remain labeled as such.

\section{A live comparison needs independent mechanisms and complete costs}

An offline controller can demonstrate request validation, abstention, tool
dispatch and preservation of graph/set semantics. It cannot measure how a live
model translates a new banking problem, how often retrieval changes its graph,
or what the workflow costs. Zero model calls means those empirical quantities
are unmeasured; it does not mean accuracy is zero or costs are free.

A credible live comparison binds the exact backend/revision and inference
settings, data-custody arrangement, actual prices and approved budget, independent
held-out mechanisms or benchmark bytes, adjudicators and rubric, and qualified
numerical baselines. The final prompts, retrieval corpus, tools, retry rules,
token/time limits and numerical tolerances should be frozen before the first
final request. Development cases already exposed to the agent cannot become
held-out by receiving different filenames.

Compare data-only, model-only, retrieval-assisted, hybrid and expert-assisted
workflows under matched information and budgets. A hybrid receives additional
information when it uses a valid domain exclusion; its improvement should not
be attributed solely to reasoning. Conversely, a data-only baseline should not
be asked to identify an effect that the observed law cannot identify. The
equal-law counterexamples in Section~\ref{v13:la-equivalence} make this boundary
testable.

Metrics should include whether the translated intervention/population is valid,
whether graph output respects equivalence and latent ambiguity, whether the
downstream effect is scientifically supported, and whether abstention occurs
for the right reason. Repeated-run disagreement is relevant even when the
average score is favorable. Fluent explanations and exact string matches to
one graph are inadequate substitutes for these criteria.

Complete costs include unsuccessful calls, retries, retrieval, tool computation,
human review and operating overhead. If an agent frequently produces a request
that needs manual correction, its low token bill can conceal high workflow cost.
Decision value should be evaluated on feasible actions and held-out outcomes
or an explicitly justified reference, with uncertainty. A stipulated arithmetic
example explains the accounting but cannot establish provider pricing, model
calibration or bank ROI.

\section{Testing semantic failures before the expensive study}

Several low-cost deterministic cases can be exercised before live qualification.
A y0 nonidentification response must not become a zero effect. A continuous
density must not be interpreted as a binary probability. A grouped allocation
contrast must not become an individual direct effect. A statsmodels internal
covariance option must not be reported as final clustered uncertainty when the
inspected implementation supplies HC0. \cite{v13_sw_y0,v13_sw_causalgps,v13_sw_inferference,v13_sw_statsmodels}

These cases test the proposed interface and evaluation rubric. They are useful
development evidence, and their exposure should be recorded. Fresh independent
mechanisms are still necessary to evaluate generalization. Keeping those roles
separate allows an agentic workflow to improve through development while
retaining a meaningful final assessment.
